%! Author = taj
%! Date = 7/22/25

% Preamble
\documentclass[11pt]{article}

% Packages
\usepackage{amsmath}
\usepackage{graphicx}
\usepackage{float}
\usepackage{upgreek}

\title{13. Trki}
\author{Taj Šobot}

% Document
\begin{document}
    \maketitle
    \flushleft

    Pri tej vaji smo preverjali veljavnost zakona o ohranitvi gibalne količine.

    \section{Meritve}

    Mase vozičkov:
    \begin{align*}
        m_1 = (0.517 \pm 0.001) kg \\
        m_2 = (0.635 \pm 0.001) kg
    \end{align*}

    V spodnjih tabelah so prikazane meritve hitrosi ob različnih fazah poskusa.
    Prožni trk:
    \begin{table}[H]
        \centering
        \begin{tabular}{|c|c|}
            \hline
            $v_1 [m/s]$ & $v_2 [m/s]$ \\
            \hline
            0.60  & 0.46 \\
            0.94  & 0.74 \\
            1.46  & 1.17 \\
            0.79  & 0.61 \\
            1.90  & 1.54 \\
            \hline
        \end{tabular}
        \caption{tabela prozni trk}
        \label{tab:1}
    \end{table}
    Neprožni trk:
    \begin{table}[H]
        \centering
        \begin{tabular}{|c|c|}
            \hline
            $v_1 [m/s]$ & $v_{12} [m/s]$ \\
            \hline
            1.09  & 0.45 \\
            2.08  & 0.93 \\
            0.68  & 0.25 \\
            1.16  & 0.51 \\
            1.60  & 0.71 \\
            \hline
        \end{tabular}
        \caption{tabela neprozni trk}
        \label{tab:2}
    \end{table}

    \section{Rezultati}
    Prožni trk:
    Preverimo da velja enačba
    \[
        m_1 v_1 = m_2 v_2
    \]
    Vstavimo za 5 trkov:
    \begin{align*}
        trk 1: (0.31 \pm 0.01) kg m/s = (0.29 \pm 0.01) kg m/s \\
        trk 2: (0.48 \pm 0.01) kg m/s = (0.47 \pm 0.01) kg m/s \\
        trk 3: (0.75 \pm 0.01) kg m/s = (0.74 \pm 0.01) kg m/s \\
        trk 4: (0.41 \pm 0.01) kg m/s = (0.38 \pm 0.01) kg m/s \\
        trk 5: (0.98 \pm 0.01) kg m/s = (0.96 \pm 0.01) kg m/s \\
    \end{align*}
    Energija, ki se je zgubila ob vsakem trku (v obliki toplote):
    \[
        \Delta W_k = \frac{1}{2}(m_2 v_2^2 - m_1 v_1^2)
    \]
    Vstavimo:
    \begin{align*}
        \Delta W_{k1} = -0.02 J \\
        \Delta W_{k2} = -0.05 J \\
        \Delta W_{k3} = -0.12 J \\
        \Delta W_{k4} = -0.05 J \\
        \Delta W_{k5} = -0.18 J \\
    \end{align*}
    Rezultati se strinjajo z teorijo.
    \\
    Isto ponovimo za neprožni trk:
    Enačba, ki more veljat:
    \[
        m_1 v_1 = (m_1 + m_2) v*2
    \]
    \begin{align*}
        trk 1: (0.51 \pm 0.01) kg m/s = (0.48 \pm 0.01) kg m/s \\
        trk 2: (1.07 \pm 0.01) kg m/s = (1.04 \pm 0.01) kg m/s \\
        trk 3: (0.33 \pm 0.01) kg m/s = (0.28 \pm 0.01) kg m/s \\
        trk 4: (0.60 \pm 0.01) kg m/s = (0.58 \pm 0.01) kg m/s \\
        trk 5: (0.82 \pm 0.01) kg m/s = (0.82 \pm 0.01) kg m/s \\
    \end{align*}
    Ter izgube energij:
    \begin{align*}
        \Delta W_{k1} = -0.19 J \\
        \Delta W_{k2} = -0.42 J \\
        \Delta W_{k3} = -0.17 J \\
        \Delta W_{k4} = -0.20 J \\
        \Delta W_{k5} = -0.21 J \\
    \end{align*}

    Opazimo, da se G ohranja (razen za to koliko izgube poberejo).
\end{document}